\chapter{Nhập Bài Không Cần Đạo Cụ}
\chapterintro{Không có đồ vẫn có cách mở lớp đẹp}{Nhiều tiết học cần khởi động ngay trong điều kiện tối giản. Khi đó, bảng, phấn, tay, mắt và chỗ ngồi là đủ.}

\section{Ba từ trên bảng}
\begin{khoidongbox}{Hoạt động khởi động}
Viết lên bảng ba từ khóa rời và mời lớp đoán bài hôm nay sẽ đi theo hướng nào.
\end{khoidongbox}
\begin{visaohieuqua}
Ba từ khóa tạo độ mở vừa đủ để lớp phải nghĩ mà không bị ngợp. Đây là cách nhập bài rất gọn và hiệu quả.
\end{visaohieuqua}
\begin{cachlam5phut}
Chọn ba từ thật sát bài. Mời 2 đến 3 dự đoán rồi nối chúng lại thành mục tiêu của tiết học.
\end{cachlam5phut}
\ngancach

\section{Bàn tay làm sơ đồ}
\begin{khoidongbox}{Hoạt động khởi động}
Dùng chính bàn tay để chia nhanh 5 ý, 5 bước, hoặc 5 dấu hiệu liên quan tới bài.
\end{khoidongbox}
\begin{visaohieuqua}
Khi cơ thể tham gia, lớp nhớ tốt hơn. Bàn tay là đạo cụ sẵn có và rất hợp với những tiết cần vào nhịp nhanh.
\end{visaohieuqua}
\begin{cachlam5phut}
Cho lớp giơ tay cùng làm, rồi chọn một ngón làm điểm mở đầu để vào nội dung chính.
\end{cachlam5phut}
\ngancach

\section{Vỗ tay chia nhịp}
\begin{khoidongbox}{Hoạt động khởi động}
Dùng một nhịp vỗ tay đơn giản để gọi lớp về cùng tốc độ trước khi bắt đầu.
\end{khoidongbox}
\begin{visaohieuqua}
Âm thanh ngắn và có nhịp kéo sự chú ý tốt hơn nhiều câu nhắc rời rạc. Nó đặc biệt hợp với lớp đông.
\end{visaohieuqua}
\begin{cachlam5phut}
Chỉ dùng trong vài giây rồi dừng hẳn. Sau nhịp vỗ, nói ngay yêu cầu đầu tiên của tiết học.
\end{cachlam5phut}
\ngancach

\section{Lấy ngay thứ trong lớp làm ví dụ}
\begin{khoidongbox}{Hoạt động khởi động}
Chỉ vào một đồ vật có sẵn như cửa sổ, bảng, quạt, ghế và hỏi nó liên quan gì đến bài hôm nay.
\end{khoidongbox}
\begin{visaohieuqua}
Ví dụ gần lớp học làm bài bớt xa. Học sinh nhìn thấy thứ trước mắt nên dễ nối ý và dễ nói hơn.
\end{visaohieuqua}
\begin{cachlam5phut}
Chọn vật thật gần, lấy 2 câu trả lời ngắn, rồi dùng luôn vật đó làm ví dụ mở đầu cho bài.
\end{cachlam5phut}
\ngancach

\section{Dùng chỗ ngồi làm dữ liệu mở bài}
\begin{khoidongbox}{Hoạt động khởi động}
Cho lớp quan sát chính cách mình đang ngồi: theo dãy, theo bàn, theo tổ, rồi đặt một câu hỏi mở từ đó.
\end{khoidongbox}
\begin{visaohieuqua}
Khai thác ngay thứ đang có giúp tiết học vào rất thật và không cần chuẩn bị thêm gì. Học sinh cũng thấy bài học gắn với không gian mình đang ở.
\end{visaohieuqua}
\begin{cachlam5phut}
Hợp với các tiết cần ví dụ về cấu trúc, phân loại, quan sát, hoặc làm việc nhóm. Chốt thật nhanh rồi vào bài.
\end{cachlam5phut}